% Shared implementation appendix. Both active papers include this file in full.
\section{Runtime and Initial Harness}
\label{app:surfaces}

The implementation described here is the repository's main branch at revision
\texttt{5ae009e9}, with both study configurations inspected on November 15,
2026. These are implementation specifications for the forthcoming
study, not experimental results. The study's eventual launch revision, frozen
configuration, split manifest, and initial/final harness hashes must accompany
its results.

\paragraph{Initialization.}
\hzero{} supplies a minimal REPL contract and generic instruction strings, a
disabled runtime policy, the default metadata formatter, empty helper and skill
libraries, and accepting answer middleware. \hzerostar{} replaces S1 with the
upstream-style system prompt, leaves S2--S5 empty, and enables the orchestrator
addendum. \texttt{H0*R}, a hand-written reference, additionally replaces
S1's recursive-call description and supplies a hand-authored S2 example. That
example distinguishes a bare completion from a child with its own REPL, splits
line-oriented input into batches of 300, requests label counts, merges counts in
Python, and retries a chunk after a parse failure. It was informed by development
observations. Optimization starts from \hzero{}, not from \texttt{H0*R}, so the
learned harness can be compared against the hand-written one rather than seeded
by it. The \texttt{initial} evaluation condition loads the exact
saved round-one harness rather than substituting a reference from today's
registry. Figure~\ref{fig:edit-surfaces} and Figure~\ref{fig:full-repl-contract}
show the sparse \hzero{} reference.

\paragraph{Execution contract.}
S1--S5 are concatenated into one system prompt; a clause referring to early turns
is guidance, not a separately timed host hook. Recursive \texttt{rlm\_query}
children below the depth limit receive their own REPL and the shared prompt.
Bare \texttt{llm\_query} calls and leaves at the recursion limit are completions,
without that REPL. Printed code output joins accumulated message history;
input-as-variable does not prevent code from printing input records. The normal
answer path uses the REPL answer dictionary, while iteration exhaustion can
invoke the runtime's final-answer fallback. Admission checks preserve specified
interfaces and caps, not a formal proof of arbitrary generated code's behavior.

% AAAI: full-width table* for the long per-surface descriptions.
\begin{table*}[t]
\centering\small
\renewcommand{\arraystretch}{1.12}
\begin{tabularx}{\textwidth}{@{}p{0.17\textwidth}X@{}}
\toprule
Surface & Editable payload, visible information, and execution scope \\
\midrule
S1: REPL contract & Full text replacement describing the environment/API; shared root/recursive-child prompt. \\
S2: Decomposition & Full text replacement guiding input inspection and decomposition; shared prompt, scoped to the current subtask. \\
S3: Execution & Full text replacement guiding parsing, calls, aggregation, and task predicates; shared prompt. \\
S4: Verification & Full text replacement guiding checks before submission; shared prompt, with no independent semantic oracle. \\
S5: Recovery & Full text replacement guiding actions after visible errors; cannot extend a terminated run or add host retries. \\
S6: Runtime policy & Dictionary with enabled flag, prompt/batch caps, depth, syntax-retry/count, and output-validation settings. Enforcement is root local except inherited depth. Width/length caps refuse work; syntax retry repeats the same prompt. \\
S7: Metadata & Root formatter function receiving bounded stdout and redacted variable types/lengths; preserves its declared output bound. It cannot inspect hidden variable values or the full call history. \\
S8: REPL helpers & One named function added/replaced in one root or child helper dictionary. Helpers execute only when called; the proposer must identify a caller and input/output contract. \\
S9: Answer middleware & Root function receiving answer text and redacted inventory; may accept transformed text or redirect with a nudge. It cannot inspect hidden records or classifications. \\
S10: Skills & One named text procedure added, replaced, or removed. A shared name/description index advertises it; its body matters only when loaded or forwarded. \\
\bottomrule
\end{tabularx}
\caption{Ten surfaces and their actual capabilities. Provider/model selection,
decoding/reasoning settings, verifiers, and experiment-owned caps are outside the
edit space. S8 spans two fields; a surface is not necessarily one builder.}
\label{tab:surfaces}
\end{table*}

\paragraph{Skills and token accounting.}
The S10 index enters the system prompt after S5. For a nonempty library the
runner installs \texttt{load\_skill(name)} in root and recursive-child REPLs.
The host merges a single-entry edit with the incumbent library, including
explicit removal. Limits are eight entries, 40 characters per name, 200 per
description, 4,000 per body, and 16,000 total. A loaded body stored only in a
variable need not be model visible. If printed or forwarded it incurs input
tokens, potentially on later turns while it remains in history. Neither loading
a body nor advertising it proves that its procedure was followed.

% Figure 1: the sparse reference harness H0 and its ten editable surfaces.
% AAAI: full-width figure* (code lines are too wide for one column); figure*
% only supports [t] placement in AAAI (stfloats is forbidden).
% Body only; the `harness' listing style, its colours, and \hzero live in the
% preamble of proposal.tex, matching how Figure_2.tex relies on the tikz styling
% defined there.
\begin{figure*}[t]
    \centering
    \begin{lstlisting}[style=harness, gobble=4]
    # shrlm/rlm_harness.py: the ten surfaces Self-Harness may edit (S1 to S10).
    
    # Condensed for visibility; the full literal appears in the appendix.
    MINIMAL_REPL_CONTRACT = """
    You are an RLM operating through a persistent Python REPL.
    Use `context`, flat or recursive LM calls, `SHOW_VARS()`, and `answer`.
    Execute one ```repl``` block per turn; print inspections and mark the
    answer ready only when complete. [Display-only condensation.]
    """
    
    def build_repl_contract() -> str:                                    # S1
        return MINIMAL_REPL_CONTRACT
    
    def build_decomposition_instruction() -> str:                        # S2
        return ("Start by probing `context` to understand what you have, "
                "then decide how the task breaks into steps.")
    
    def build_execution_instruction() -> str:                            # S3
        return ("On each turn, run one block of code, print what you need "
                "to see, and build the result up in REPL variables.")
    
    def build_verification_instruction() -> str:                         # S4
        return 'Check your candidate answer before setting `answer["ready"] = True`.'
    
    def build_recovery_instruction() -> str:                             # S5
        return ("If a sub-call errors or returns something you cannot use, "
                "adjust the approach and try again.")
    
    def build_runtime_policy() -> dict[str, Any]:                        # S6
        return {
            "enabled": False,
            "max_prompt_chars": None,
            "max_batch_width": None,
            "max_depth": None,
            "retry_on_syntax_error": None,
            "max_retries": None,
            "validate_sub_output": None,
        }
    
    def build_metadata(                                                  # S7
        stdout: str,
        repl_inventory: dict[str, tuple[str, int]],
        max_character_length: int = DEFAULT_MAX_CHARACTER_LENGTH,
    ) -> str:
        return default_metadata_builder(stdout, repl_inventory,
                                        max_character_length)
    build_metadata.declared_bound = DEFAULT_MAX_CHARACTER_LENGTH        # S7
    
    def build_repl_helpers() -> dict[str, Any]:                          # S8
        return {}
    
    def build_sub_repl_helpers() -> dict[str, Any]:                      # S8
        return {}
    
    def build_answer_middleware() -> AnswerMiddlewareFn:                 # S9
        return accept_answer
    
    @dataclass(frozen=True)
    class SkillEntry:                                                    # S10
        name: str          # REPL-safe identifier; index label
        description: str   # one line: when to consult it; enters the prompt
        body: str          # ordered steps; served verbatim by `load_skill(name)`
    
    def build_skills() -> list[SkillEntry]:                              # S10
        return []
    \end{lstlisting}
    \caption{Sparse reference harness \hzero{} and its ten editable surfaces,
    each a declared field or group of fields (S8 has two builders). The configured
    study starts from \texttt{H0*R}, described in Appendix~\ref{app:surfaces}.
    S1 is condensed for display; Figure~\ref{fig:full-repl-contract} gives the
    full \hzero{} S1 template. S10 advertises a name/description index and loads
    bodies on demand; token cost depends on exposure in conversation history or
    forwarding to children. Each admitted edit changes one surface, and a round
    may compose several edits on distinct surfaces into one measured candidate.}
    \label{fig:edit-surfaces}
\end{figure*}

% Appendix figure: exact, unabridged sparse H0 S1 template (not H0*R).
% Figure 1 intentionally contains a display-only condensation of this prompt.
% AAAI: full-width figure* (the prompt lines are too long for one column);
% figure* only supports [t] placement in AAAI (stfloats is forbidden).
\begin{figure*}[t]
    \centering
    \begin{lstlisting}[style=harness, gobble=4]
    MINIMAL_REPL_CONTRACT = textwrap.dedent(
        """\
        You are a Recursive Language Model (RLM): a language model whose prompt-related information lives in a Python REPL that you drive turn by turn. You will be queried turn-by-turn until you have an answer to the query.

        Available in the REPL:
        - `context`: the information related to the prompt (typically `str` or `list[str]`).
        - `llm_query(prompt: str, model: str | None = None) -> str`: a single sub-LLM completion.
        - `llm_query_batched(prompts: list[str], model=None) -> list[str]`: several sub-LLM completions in parallel; same order out as in.
        - `rlm_query(prompt, model=None)` / `rlm_query_batched(prompts, model=None)`: recursive RLM sub-calls. Fall back to `llm_query` / `llm_query_batched` when recursion is disabled.
        - `SHOW_VARS() -> str`: list every variable currently in the REPL.
        - `answer`: dict initialized to `{{"content": "", "ready": False}}`. To submit, set `answer["content"]` to the final answer and `answer["ready"] = True` inside a ```repl``` block.
        {custom_tools_section}

        Write code in ```repl``` blocks, one block per turn; the REPL persists across turns. The REPL is NOT a Jupyter cell - only `print(...)` output (stdout) is shown back to you between turns; a bare expression on the last line is silently discarded. REPL outputs over ~20K characters are truncated."""
    )
    \end{lstlisting}
    \caption{The complete sparse \hzero{} S1 source template, before placeholder
    formatting. Figure~\ref{fig:edit-surfaces} uses a display-only condensation.
    This is a reference contract, not the full initial prompt of the configured
    \texttt{H0*R} study; its upstream-style prompt, recursive-call change, S2
    example, and orchestrator addendum are specified in
    Appendix~\ref{app:surfaces}.}
    \label{fig:full-repl-contract}
\end{figure*}

\section{Mining, Diagnosis, and Evidence Selection}
\label{app:mining}

Final-answer verification supplies pass/fail and an outcome category. Eligible
failed runs are sent to a fixed-model attributor; passing runs and configured
environment-owned exclusions are not sent through failure attribution. Excluded,
unattributed, and rejected records remain visible in accounting. Each diagnosis
records an inferred mechanism, location and causal status, supporting operation
references, and uncertainty. Structural absence of recursion can be observed
directly; semantic correctness and root-versus-child causality are not implied
by final-answer scoring.

The attributor must distinguish skipped records, wrong or uncertain labels, and
incorrect aggregation or predicates. Missing answer elements alone do not prove
input loss. It must account for observed checks and repairs, and distinguish
missing evidence from evidence that an action was omitted. For
\texttt{incomplete\_coverage}, a structured basis records status, input scope,
loss observation, and counterevidence. Statuses \texttt{not\_established} and
\texttt{contradicted} normalize to \texttt{OTHER}/\texttt{UNATTRIBUTED}, rather
than authorizing a coverage edit. Host acceptance verifies schema, vocabulary,
field bounds, and resolvable references, not the truth of a causal explanation.
Attribution response attempts default to three; transport retries are separate.

\paragraph{Clustering and routing.}
Records cluster by exact agreement on
% AAAI: split over two lines so it fits the single column.
\[
\begin{aligned}
\phi(r)=(&\text{outcome cause},\ \text{recorded failing level},\\
          &\text{causal status},\ \text{mechanism}).
\end{aligned}
\]
The signature mixes observed outcomes with model-inferred fields; it is not a
verified causal label. Patterns rank first by distinct-instance support, then
actionability, then a deterministic signature tie-break. Attempt support is
recorded separately. Minimum support two is a flag, not automatic deletion.
Three representatives and a fixed actionability heuristic guide reading order:
weights are 0.4 causal status, 0.3 grounding availability, 0.2 known mechanism,
and 0.1 normalized-description homogeneity. The last term uses distinct text
strings, not semantic embeddings; the score is not a calibrated success
probability.

\begin{table}[t]
\centering\small
\begin{tabularx}{\linewidth}{@{}Xl@{}}
\toprule
Mechanism & Eligible surfaces (primary first) \\
\midrule
REPL contract misuse & S1, S9, S10, S4 \\
Task objective misinterpretation & S2, S3, S4, S10 \\
Whole-input subcall collapse & S2, S3, S10, S6 \\
Incomplete coverage & S2, S4, S3, S10 \\
Redundant decomposition & S2, S3, S10 \\
Misaligned unit & S2, S10 \\
Depth degradation & S3, S6, S2 \\
Insufficient recursion & S3, S2, S6 \\
Bare completion substituted for recursion & S3, S1, S10 \\
Premature termination & S4, S9, S3 \\
Skipped verification & S4, S9, S10 \\
Swallowed subcall error & S5, S3, S6 \\
Iteration budget exhaustion & S6, S3, S1, S10, S5$^{\dagger}$ \\
Unparsed child output & S7, S3, S5, S8$^{\dagger}$ \\
REPL execution fault & S8, S10, S3, S1, S5$^{\dagger}$ \\
Answer serialization failure & S9, S4, S10 \\
Lossy aggregation & S3, S4, S8$^{\dagger}$ \\
Unconsulted procedure & S10, S3, S2 \\
Other & S1--S10 \\
\bottomrule
\end{tabularx}
\caption{Human-defined, many-to-many routing (taxonomy 3.4.0).
An \texttt{UNATTRIBUTED} diagnosis has no eligible route, including when its
mechanism is Other. $\dagger$ denotes additional operation-evidence requirements.}
\label{tab:routes}
\end{table}

Conditional S8 eligibility requires a complete cited operation; unparsed child
output additionally requires child-response evidence. Conditional S5 eligibility
requires an observed error and a subsequent complete operation. These are
structural opportunities to intervene, not proof of successful recovery or an
effective proposed fix. Aggregation errors route to S3/S4 or supported S8;
S9 remains limited to answer-visible properties.

\paragraph{Bounded evidence.}
The attribution digest defaults to 12,000 characters. Proposal evidence defaults
to 32,000 characters, separately from the 12,000-character history budget.
These are component budgets, not a total prompt cap. A compact pattern inventory
is supplemented by up to $\min(k,4)$ expanded patterns, preferring distinct
actionable mechanisms before additional signatures/context. Inventory-only
patterns do not authorize proposals. Evidence follows cited child calls and
relevant parsing, merging, or filtering operations, with nearby input definitions,
consumers, and corrections. Whole code operations are retained where they fit;
required operations that cannot fit are omitted with scope markers rather than
presented as complete fragments. Outputs and prose may be shortened. A passing
contrast sharing operation names can be included when available. This improves
context for comparison but guarantees neither full causal context nor balanced
surface coverage.

\section{Proposal, Materialization, and History}
\label{app:proposal}

One response under \texttt{proposal-selection/v2} selects evidence-supported
pattern/surface interventions before supplying replacement payloads; these are
not separate selection and generation calls. Each admitted candidate occupies
one exact (pattern, surface) pair: the same pattern may be attacked on several
surfaces, and the same surface may carry edits for different patterns, but a
pair cannot be claimed twice. Rejected entries reserve nothing. At most $k$
candidates are admitted, with fewer or zero allowed. Every admitted candidate is
materialized as its own complete harness, the incumbent with one surface
changed; candidates are never composed with each other. Three short explanation
fields identify incumbent behavior, unresolved observed failure, and the change
at a specific operation, input, or trigger. Guidance asks whether the same
demonstrated failure could persist if the edit were followed perfectly, and
requires the change to be stated as a forcing rule, a concrete trigger the
model will actually meet on a real input, rather than as softer or more
emphatic phrasing. Repeating an existing instruction more emphatically, moving
it to another surface, or repairing an already-recovered fault is insufficient
without a substantive behavioral difference. Replacing a misleading example can
be the smallest effective change.

The general prompt asks which information must survive each operation and which
task conditions must the final computation enforce. It does not prescribe an
OOLONG-specific record-ID algorithm. Environment-provided API and answer
contracts remain authoritative. These instructions guide model judgment; they
are not deterministic tests of semantic novelty or correctness.

\paragraph{Repair and preflight.}
Response/schema attempts are bounded (default maximum eight), distinct from
task validations. When failed-member repair is invoked there is one repair
response, restricted to failed original patterns. It may select another eligible
unoccupied surface or withdraw an edit while retaining valid siblings unchanged;
it may not invent new patterns. Collision repair must choose a contender or
withdraw it, still enforcing one admitted candidate per (pattern, surface) pair.

For text under \texttt{literal-text/v1}, the model supplies literal instructions;
the host encodes formatting braces exactly once and restores the designated
custom-tools placeholder. S1--S5 replace full text, S6 replaces a policy
dictionary, S7/S9 supply functions with required interfaces, S8 edits one named
helper in one dictionary, and S10 edits one named library entry. Preflight checks
serialization and base hashes, an actual single-surface change, reserved
interfaces, caps, formatting, and selected smoke/answer fixtures. Environment
answer fixtures are synthetic contract examples, not validation tasks. Literal
no-op rejection does not detect every behaviorally redundant instruction.
Saved artifacts distinguish literal proposal text, encoded source/template, and
runtime rendering.

\paragraph{History and qualified progress.}
The archive retains attempts, admission/rejection reasons, screening ranks,
measured outcomes, and revision links. The bounded prompt summary prioritizes
recent rounds, the incumbent's promotion, relevant measured predecessors, and an
older promising direction when it fits. Mechanisms already attempted on more
than one surface are listed separately, so the proposer sees that moving the
same idea to a new surface has been tried. Revisiting an intervention requires a
resolvable predecessor reference and a stated substantive change to the operation
or joint hypothesis. Validation contributes aggregate scores, costs, sanitized
failure summaries, and observed activation; its raw traces and gold/answer dumps
are not forwarded to the proposer.

A rejected measured candidate can be marked potentially promising when a
comparable verifier-defined dense metric improves in its declared direction.
Definitions specify identity/version, direction, aggregation, precision, and
terminal-outcome values. Comparison also requires compatible verifier/protocol/
repetition settings and matching instance-attempt sets. Missing, unsupported, or
incomparable metrics remain \texttt{not\_assessed}. The label is qualified
evidence, not promotion or statistical significance; a simultaneous exact-pass
regression remains recorded. An identical rejected harness already measured
under the same incumbent can be skipped before new validation.

\paragraph{Observed behavior.}
Instrumentation reports root S6 syntax retries and S9 redirects, and S10 loads
across the traced call tree, with observation coverage. Missing coverage is
unassessed. The intended-behavior field remains \texttt{not\_assessed}: activation
does not establish compliance or task correctness. There is no general semantic
compliance detector for S1--S5, S7, or S8.

\section{Screening, Validation, and Stopping}
\label{app:validation}
\label{app:algorithm}

Let $H$ be the incumbent and $C_1,\dots,C_{k'}$ ($k'\le k$) the admitted
candidates, each $H$ with one surface replaced. Let $V$ be the held-out
validation tasks, stratified by optimization context length
$\ell\in\mathcal{L}$, and $V_s\subset V$ a small fixed screening subset. For
OOLONG-Pairs, $V_s$ is curated from held-out task IDs of known difficulty and
spans every length (one, two, and three instances at 8K, 16K, and 32K). For
OBLIQ it holds the hardest and a median-scoring validation query per pool size.
Other environments use a generic small subset. Validation tasks are never mined.

\paragraph{Promotion metric and rule.}
The configured metric is either the verifier-pass count or a verifier-owned
primary quality. Both study profiles use primary quality: pair-set F1 for
OOLONG-Pairs and NDCG@10 for OBLIQ, written $F$ below. For harness $X$ let
$\bar F_\ell(X)$ be its mean primary quality over the attempts in stratum $\ell$, and
$\Delta_\ell=\bar F_\ell(C)-\bar F_\ell(H)$. A candidate $C$ is accepted when
\begin{equation}
\Delta_\ell\geq-\tau_\ell\ \ \forall\ell\in\mathcal{L},\qquad
\bar\Delta=\tfrac{1}{|\mathcal{L}|}\textstyle\sum_{\ell}\Delta_\ell\geq\tau_{\mathrm{imp}},
\label{eq:promotion}
\end{equation}
with equality tolerated up to floating-point precision. The per-length margins
$\tau_\ell$ are calibrated by a nested bootstrap over the incumbent's own
held-out scores at the 95th percentile, so any one length can veto a candidate
whose macro mean improves. Under the pass-count metric there is no length axis,
and a single $\tau_{\mathrm{reg}}$ bounds the pass-count change. For each
configured resource metric $q$, the candidate mean must also lie in the
inclusive band $a_q\bar q(H)\leq\bar q(C)\leq b_q\bar q(H)$, with unbounded
endpoints permitted. Cost-band enforcement can be disabled, in which case the
band is still recorded. An optional list of candidate-owned terminal failure
causes may not become more frequent than under $H$. Per-profile values are in
Table~\ref{tab:study-settings}; neither profile constrains sub-calls. Under the
OOLONG-Pairs setting $\tau_{\mathrm{imp}}=0$, a candidate with $\bar\Delta=0$
that clears every $\tau_\ell$ promotes. This is a selection rule, not a significance
test or a guarantee of per-task improvement.

\paragraph{Screening tournament.}
If $k'\ge2$, $H$ and every $C_i$ are evaluated independently on $V_s$ with
$v_s<v$ attempts per task; the incumbent is run fresh on exactly $V_s$ so that
every comparison uses the same tasks and attempts. Each candidate is scored
against $H$ with Equation~\ref{eq:promotion} and ranked by $\bar\Delta$.
Candidates that were never scored (for example, over budget) sink below every
scored one, and ties break deterministically by candidate identifier. Screening
verdicts are written to a separate screening ledger and are never promotion
decisions. If $k'\le1$, the tournament is skipped and the single candidate (if
any) proceeds directly.

\paragraph{Re-attribution.}
If every screened candidate fails Equation~\ref{eq:promotion} on $V_s$, the
already-executed mining runs are attributed a second time into a separate
evidence bundle, with no new task rollouts. This targets attribution noise:
independent re-attributions of the same digest often disagree. The proposer
then produces a fresh candidate pool from that bundle, which is screened in the
same way. The retry happens at most once per round. If it yields no candidate,
the original ranking stands. If its pool is also screened out, its own ranking
is used.

\paragraph{Full validation.}
Only the top-ranked candidate $C^\star$ is evaluated on all of $V$ with $v$
attempts, alongside $H$. Its $v_s$ screening attempts on $V_s$ are copied into
this evaluation and only the remaining instance-attempts run. If $H$ has not
changed since an earlier round, its completed held-out runs are reused rather
than repeated. Screening runs never seed an incumbent baseline, because they
cover only $V_s$. Under F1 promotion with early stopping enabled, $H$ is
evaluated first and $C^\star$'s tasks are ordered by length. Once all
of a length's runs are complete and $\Delta_\ell<-\tau_\ell$, the remaining
runs are skipped and $C^\star$ is rejected. Early stopping changes cost, not
the rule: a candidate that completes is judged by
Equation~\ref{eq:promotion}. If $C^\star$ is rejected, the other screened
candidates are retired rather than validated in turn.

\paragraph{Refinement.}
A rejected $C^\star$ may still improve F1 on runs that produced a measurable
answer while failing more often with a terminal cause such as malformed
output. Refinement is considered when its F1 conditional on a measurable answer
exceeds $H$'s by at least $\delta_{\mathrm{ref}}$ on the same length strata
($\delta_{\mathrm{ref}}=0.02$). The configured terminal-cause failure rate is
recorded but does not gate admission. One repair generation then mines
$C^\star$ on the held-in tasks and attributes the failures. The proposer is
restricted to patterns whose verifier cause is in the configured set
(\texttt{wrong\_format}) and to a single edit applied to $C^\star$. The
resulting descendant is validated against the \emph{original} $H$ on the same
$V$ under Equation~\ref{eq:promotion}, reusing $H$'s held-out runs. The branch
has at most one generation and never recurses. Because it spends a full
mining-and-validation leg, the implementation waits for an explicit recorded
confirmation before running it. The decision is persisted either way.

\begin{quote}\small
\textbf{Inputs:} initial harness $H$, mining tasks $M$, validation tasks $V$ and
screening subset $V_s$, repetitions $m,v,v_s$, candidate ceiling $k$, round
ceiling $T$, patience $p$, and promotion settings. Initialize history and
no-promotion counter $z=0$.
\begin{enumerate}
\item For each round up to $T$, execute or resume $H$ on $M$ with $m$ attempts;
record final outcomes, traces, and available observations.
\item Diagnose eligible failures, cluster records, and select bounded evidence
and history. Propose, admit, and if needed repair up to $k$ single-surface
candidates on distinct (pattern, surface) pairs; preflight each one.
\item If at least two candidates load, screen $H$ and every candidate on $V_s$
with $v_s$ attempts and rank them. If all are screened out, re-attribute the
same mining runs once, propose a fresh pool, and re-screen.
\item Fully evaluate the top-ranked $C^\star$ (or the only candidate) and $H$
on $V$ with $v$ attempts, reusing $C^\star$'s screening runs and any stored runs
of an unchanged $H$, with optional per-length early stopping. An unscorable
baseline follows the experiment error path. If Equation~\ref{eq:promotion} and
the resource bands hold, set $H=C^\star$.
\item Otherwise, if refinement is enabled and admitted, run one confirmed repair
generation from $C^\star$ and apply the same rule to its descendant against $H$.
\item Persist screening and promotion ledgers, or the reason no comparison
occurred; update history. Set $z=0$ on promotion, otherwise $z=z+1$. Stop when
$z\geq p$.
\item Freeze the last incumbent; run separately specified final-evaluation
conditions. Do not retrospectively select a harness using test scores.
\end{enumerate}
\end{quote}

\paragraph{Failures, denominators, and resume.}
Candidate-owned execution errors can become failed attempts without terminating
the whole experiment. Configuration/operational/provider errors, cancellation,
memory or I/O failures, and observation-persistence failures may still propagate.
Complete comparisons require matching planned attempts; a partial baseline cannot
anchor promotion. The final-evaluation driver can retain completed sibling
results while reporting an incomplete or over-budget condition.

\begin{table}[t]
\centering\small
\begin{tabularx}{\linewidth}{@{}p{0.27\linewidth}X@{}}
\toprule
Recorded status & Treatment \\
\midrule
Completed correct/incorrect & Enters attempt denominator; only a verifier pass counts as success. \\
Completed execution/resource/format/filter failure & Enters denominator as a failure; dense terminal values follow the verifier definition. Partial usage/trace coverage remains labeled. \\
Missing/canceled attempt & Report separately, not as a fabricated measured zero. Incomplete validation cannot silently become a complete comparison. \\
Unscorable baseline & Persist evidence and stop/raise through the experiment error path; this is not a candidate rejection. \\
Resumed/cached completion & Reuse only with matching identity/contracts; not a new independent attempt or new spend. \\
\bottomrule
\end{tabularx}
\caption{Outcome accounting for optimization and final evaluation.}
\end{table}

\section{Experimental Protocol and Provenance}
\label{app:protocol}

\paragraph{Profiles.}
Optimization runs from
\path{configs/experiment_oolong_pairs_blended_openrouter_DeepSeekV4Flash.toml} and
\path{configs/experiment_obliq_bench_math_openrouter_DeepSeekV4Flash_blended_h0init.toml}
via \path{examples/run_experiment.py}. Final evaluation runs from
\path{examples/run_evaluation.py} on each experiment's frozen splits, with conditions
\texttt{b1} (\hzero{}), \texttt{h0\_star} (\hzerostar{}), \texttt{lambda\_rlm}, and
\texttt{sh\_rlm}. The configuration identity is hashed on first use, and any later
change to an identity-covered setting refuses to resume the experiment.

\paragraph{OOLONG-Pairs data and scoring.}
The dataset is \texttt{oolongbench/oolong-synth}, upstream validation partition,
\texttt{trec\_coarse}, revision
\path{f0d59eaf0febf130664cfceb710436c8e3216b2b}. For each nominal window (8K,
16K, 32K, and the test-only 262K), the loader selects the first two distinct
record collections, with a 50,000-row scan ceiling, and pairs each with task
IDs 1--20, giving 40 instances per window. Labels are removed from
model-visible records; gold user-pair sets are computed from the labeled
records and task predicates. Within each optimization window, roles are 8
mining, 8 validation, and 4 test instances. A further 8 are reserved for a
refinement promotion split that the current implementation does not use,
because refinement validates on the ordinary validation set. Role assignment
follows a curated task order rather than the raw seeded shuffle, so that
validation includes known-hard task IDs (11 and 18) as well as easier ones. All
40 262K instances are test-only.

The scorer extracts integer pairs and compares sets. Global pair order and
duplicate pairs do not affect set equality, but order within a pair is preserved:
the lower user ID must come first. A bracketed list or newline-separated pairs
are accepted. The explicit empty marker is ``No valid pairs found.''; literal
\texttt{[]} is not an empty answer. No parseable pair or marker is a format
failure. A correctly empty prediction against an empty gold set passes and has
precision/recall/F1 equal to one. These extraction and empty-set distinctions
are adaptations of the example scorer, not an unmodified upstream verifier.

Different predicates over the same record collection ask different aggregation
questions over identical records, and each window has only two collections. Test
results at the optimization windows therefore measure held-out task instances,
not unseen records. The 262K test uses larger, distinct collections that share
question texts with the shorter ones; it measures length generalization, not
wholly unseen source text. Stripping labels and adding task text changes actual
prompt lengths, so evaluation reports measured prompt sizes separately from the
nominal window labels.

\paragraph{OBLIQ-Bench Math data and scoring.}
The dataset is \texttt{dianetc/OBLIQ-Bench}, \texttt{analogues/math} subset,
revision \path{4ebee29f68ceeb62ca00bd73d5478acdd3cd7764}: 3,508 corpus problems,
151 queries, and binary relevance judgments. Each instance is a query, a
candidate pool, and instructions to end with a ranked list of at most ten
bracketed corpus IDs. The pool holds the query's gold problems plus random
negatives up to 200, 500, or 1,000 problems, sampled with a seed per query;
the query's own self-match IDs are excluded. Each pool size draws a
\emph{disjoint} set of queries rather than the same queries at three
difficulties. Test queries come from a frozen 40-query split used in earlier
full-corpus comparisons, and mining and validation queries come from the
remaining 111. Each pool size receives 8 mining, 8 validation, and 4 test
queries. The screening subset holds the hardest and a median-scoring validation
query per pool size, chosen from the round-one incumbent's validation scores.

The verifier takes the last \texttt{RANKED:} list in the answer, keeps the first
ten distinct IDs, and computes NDCG@10. A missing list is a format failure, and
an explicit empty list is \texttt{no\_answer}. NDCG@10 of at least 0.9 counts as
a pass. Below that, the failure is labeled incomplete, spurious, or mixed by
whether gold IDs are missing, extra IDs are present, or both. Because all
judgments share one relevance grade, a prediction equal to the gold set scores
1.0 regardless of order.

\paragraph{Model, caps, and concurrency.}
All three roles use \texttt{deepseek/deepseek-v4-flash} through OpenRouter with
temperature 0, top-$p$ 0.95, and an output-token cap of 8,192. The provider's
default reasoning mode is enabled, without guaranteeing returned reasoning.
Each run is capped at depth three, 3,600 seconds, and \$2. OOLONG-Pairs allows 45
iterations plus a three-iteration, 600-second finalization window in which only
answer submission remains available. OBLIQ allows 30 iterations. Each validation
or screening subject, including the incumbent, has a \$100 cumulative budget.
Concurrency settings (mining, screening, and validation workers) do not change
any repetition count. Final evaluation uses three attempts per task per
condition.

The $\lambda$-RLM condition is identified by its method kind, reconstruction
version, upstream revision, and source hash; it is not an unmodified released
implementation for either environment. Model/deployment, rendered prompts,
source/configuration snapshots, split IDs/hashes, exact edit payloads, harness
lineage, and analysis versions belong in the study artifacts. Freezing these
before a run does not by itself constitute preregistration. Previously inspected
development or test material must retain its exposure history if later used to
support a result.

\section{Measurement, Cost, and Saved Observations}
\label{app:measurement}
\label{sec:cost}

\paragraph{Quality and surface activity.}
Exact pass is set equality, not rounded F1 equal to one. OOLONG-Pairs stores
per-attempt precision, recall, and F1 to three decimals; reported mean F1 averages
recorded attempt values, including verifier-declared zeros for completed runtime,
resource, format, and content-filter failures. It is macro averaging across
attempts, not pooled pair-level micro-F1. Missing/extra counts are reported only
where available, with coverage, rather than filling absent diagnostics with zero.
OBLIQ-Bench Math stores per-attempt NDCG@10 to three decimals, with format and
empty-answer failures scored zero; its pass rule is NDCG@10 at least 0.9.
Other environments supply their own metric definitions or remain unassessed.

Surface activity distinguishes submitted, admitted, evaluated, and promoted
edits, and records each candidate's screening rank separately from its
promotion decision. Because candidates are never merged, a promoted edit is the
whole measured difference from its incumbent; screened-out and retired
candidates carry only screening-subset estimates. Per-round quality plots use the freshly measured
candidate on promotion and freshly measured incumbent otherwise; a round without
a measurement is missing, not a zero or a carried-forward score. Environment
labels come from the configuration. Trace summaries can report depth, call
counts, and the implemented whole-input-subcall proxy: maximum descendant prompt
characters divided by root prompt characters exceeds 0.8. Its denominator is
walkable runs, passing and failing, with unwalkable runs separate; it is neither
a token ratio nor proof of meaningless decomposition.

\paragraph{Cost.}
An ordinary completed round with $k'\ge2$ loaded candidates schedules at most
\begin{equation}
N_{\mathrm{round}}=m n_{\mathrm{in}}
+(k'+1)\,v_s n_s
+\underbrace{v n_{\mathrm{ho}}}_{H}
+\underbrace{(v n_{\mathrm{ho}}-v_s n_s)}_{C^\star}
\label{eq:round-cost}
\end{equation}
task attempts, where $n_s=|V_s|$. The incumbent term is zero when an unchanged
incumbent's held-out runs are reused, and per-length early stopping can shorten
the $C^\star$ term. With $k'\le1$ there is no screening term and the candidate
runs all $v n_{\mathrm{ho}}$ attempts. A re-attribution adds one more screening
pool, at most $(k+1)v_s n_s$. An admitted, confirmed refinement adds
$m n_{\mathrm{in}}+v n_{\mathrm{ho}}$, with the incumbent's runs reused.
Both blended profiles have $n_{\mathrm{in}}=n_{\mathrm{ho}}=24$ across three
sizes, $m=2$, $v=3$, $k=4$, $v_s=1$, and $n_s=6$, so a full round schedules
$48+30+72+66=216$ attempts, or 144 when the incumbent is reused. OBLIQ disables
early stopping, re-attribution, and refinement, so its rounds have no further
terms. No admitted candidate means no screening or validation leg. Final
evaluation runs three attempts per task for each of the four conditions. With
12 short and 40 long OOLONG-Pairs test tasks that is $3\cdot52\cdot4=624$
attempts, and with 12 OBLIQ test queries it is $3\cdot12\cdot4=144$. These are
scheduling counts, not measured results or provider-call counts. Root attempts
may issue many recursive/model requests. Attribution, proposal, repairs, and
operational retries also consume resources.

The spend breaker uses the per-response cost OpenRouter returns. Configured
rates of \$0.009 per million input tokens and \$1.28 per million output tokens,
OpenRouter's listed price when the profiles were written, feed only cost
projections; they are not a current market quote. Actual usage and its coverage will be reported separately from scheduling
projections. Task usage is not added again to a stage total
that already includes it. Breaker reservations are not billed spend; interrupted
or unreported requests can leave a lower bound, and the per-run budget is not a
guarantee of complete invoice accounting.

\paragraph{Provider-returned observations.}
Instrumented execution, mining, validation, final evaluation, attribution,
proposal/repair, recursive calls, compaction, and fallback paths save returned
response/reasoning observations separately from executable code and final
answers. References carry call IDs, paths, and hashes. Run traces point to
\path{<run_id>.llm_calls/<call_id>/}; attribution and proposal artifacts point to
their stage-specific \path{llm_calls/} directories. Started, response, and finished
events distinguish failed or interrupted calls. Cache hits retain source
provenance and do not add spend.

Availability is explicit: \texttt{returned}, \texttt{opaque\_only},
\texttt{not\_returned}, \texttt{unsupported\_client},
\texttt{legacy\_unavailable}, or \texttt{no\_response}. Returned content can be
text or a summary; this is not guaranteed access to full private reasoning.
Unknown reasoning-token counts stay unknown and are not added to reported output
usage again. Legacy traces cannot be backfilled with unreturned reasoning.
Observation artifacts are not automatically fed into attribution digests,
proposal evidence, or history. Ordinary assistant prose in a digest is a
different source. Observation persistence errors remain visible failures.

% AAAI: full-width table* (the version strings are too long for one column).
\begin{table*}[t]
\centering\small
\begin{tabularx}{\textwidth}{@{}Xl@{}}
\toprule
Contract & Version \\
\midrule
Taxonomy / capabilities & 3.4.0 / \texttt{surface-capabilities/v2} \\
Attribution prompt / validator / digest & 1.6.0 / 1.2.0 / 1.8.0 \\
Proposal prompt / validator & 4.3.0 / 4.2.0 \\
Proposal response / text & \texttt{proposal-selection/v2} / \texttt{literal-text/v1} \\
Evidence selector / diagnostic history & 4.3.0 / 2.1.0 \\
Compact history & \texttt{proposal-history/v3} \\
Validation / summary & \texttt{heldout-batch/v2} / \texttt{shrlm-validation-summary/v3} \\
Behavior / provider observations & \texttt{behavior-observations/v1} / \texttt{rlm-llm-observation/v1} \\
\bottomrule
\end{tabularx}
\caption{Contracts at the documented revision; archive the versions actually
used by the forthcoming run with its results.}
\end{table*}

\section{Optional Follow-up Studies}
\label{app:extensions}
\label{app:finetune}
\label{app:alignment}

Cross-environment transfer would require a separately specified target protocol
and completed evaluation. A weight-training comparison would require an available
trainable checkpoint, justified backbone matching, and measured compute; no
fine-tuning condition or hardware allocation is claimed for this hosted study.
Because each promotion measures one single-surface edit against its incumbent,
per-round outcomes already attribute measured gains to individual edits. An optional future
leave-one-surface-out study would revert a final surface to its starting value
and rerun the same evaluation; it would measure that final-state counterfactual,
not the marginal contribution of every historical edit overwritten along the way.
No such ablation has been completed here.

An optional safety probe could compare frozen harness states for changes in
refusal and harmful-answer behavior, motivated by work on
misevolution~\citep{shao2025misevolution}. The task gate does not measure these
outcomes, so improved benchmark behavior would not establish safe deployment.
Any probe would need its own data, scoring, and budget. Textual reasoning
instructions are editable; provider reasoning settings are experiment owned.
